\documentclass[11pt,a4paper]{article}

\usepackage[utf8]{inputenc}
\usepackage[T1]{fontenc}
\usepackage[polish]{babel}
\usepackage{geometry}
\usepackage{amsmath}
\usepackage{amssymb}
\usepackage{graphicx}
\usepackage{booktabs}
\usepackage{float}
\usepackage{tikz}
\usepackage{pgfplots}
\usetikzlibrary{shapes.geometric, arrows.meta, positioning}
\pgfplotsset{compat=1.18}

\geometry{margin=2.5cm}

\title{Inteligencja Obliczeniowa i jej Zastosowania \\ Laboratorium 1 - sprawozdanie}
\author{Jakub Jaśków}
\date{\today}

\begin{document}

\maketitle

\section{Opis analizowanego systemu}

Analizowanym systemem jest instalacja przeciwpożarowa uruchamiana na żądanie w budynku użyteczności publicznej. System składa się z czujników dymu, centrali sterującej oraz instalacji gaśniczej (np. tryskaczowej).

System pracuje w trybie \textit{on-demand}, ponieważ jego działanie wymagane jest jedynie w sytuacji wystąpienia pożaru. W normalnych warunkach system pozostaje w stanie gotowości.

Zdarzeniem niepożądanym analizowanym w pracy jest:

\[
T = \text{brak uruchomienia instalacji gaśniczej podczas pożaru}
\]

Zdarzenie to może prowadzić do strat materialnych, zagrożenia zdrowia i życia ludzi oraz niespełnienia wymagań bezpieczeństwa.

\begin{figure}[H]
\centering
\begin{tikzpicture}[
    block/.style={rectangle, draw, rounded corners, minimum width=3cm, minimum height=1cm, align=center},
    >=Latex,
    node distance=2cm
]

\node[block] (sensor) {Czujnik dymu};
\node[block, right=of sensor] (control) {Centrala\\sterująca};
\node[block, right=of control] (fire) {Instalacja\\gaśnicza};

\draw[->] (sensor) -- (control);
\draw[->] (control) -- (fire);

\end{tikzpicture}
\caption{Schemat działania analizowanego systemu przeciwpożarowego uruchamianego na żądanie.}
\label{fig:system}
\end{figure}

\section{Rozmyta skala prawdopodobieństwa}

Do opisu prawdopodobieństw zdarzeń bazowych zastosowano zmienną lingwistyczną \textit{Probability}. Każdej wartości lingwistycznej przypisano trapezową funkcję przynależności.

\begin{center}
\begin{tabular}{ccc}
\toprule
Wartość lingwistyczna & Opis & Parametry trapezu $(a,b,c,d)$ \\
\midrule
VS & bardzo małe & $(0,0,0.002,0.01)$ \\
S  & małe & $(0.005,0.01,0.03,0.05)$ \\
M  & średnie & $(0.03,0.06,0.10,0.15)$ \\
B  & duże & $(0.10,0.15,0.20,0.30)$ \\
VB & bardzo duże & $(0.25,0.30,0.40,0.50)$ \\
\bottomrule
\end{tabular}
\end{center}

W dalszej części sprawozdania każdą liczbę rozmytą zapisano w postaci trapezowej:
\[
\tilde{X}=(a,b,c,d)
\]

gdzie:
\begin{itemize}
\item $[a,d]$ jest nośnikiem liczby rozmytej,
\item $[b,c]$ jest rdzeniem liczby rozmytej.
\end{itemize}

\begin{figure}[H]
\centering
\begin{tikzpicture}
\begin{axis}[
    width=15cm,
    height=8cm,
    xmin=0, xmax=0.5,
    ymin=0, ymax=1.1,
    xlabel={Prawdopodobieństwo},
    ylabel={Stopień przynależności},
    grid=both,
    legend style={at={(0.5,-0.2)}, anchor=north, legend columns=5}
]

\addplot[thick] coordinates {(0,0) (0,1) (0.002,1) (0.01,0)};
\addlegendentry{VS}

\addplot[thick] coordinates {(0.005,0) (0.01,1) (0.03,1) (0.05,0)};
\addlegendentry{S}

\addplot[thick] coordinates {(0.03,0) (0.06,1) (0.10,1) (0.15,0)};
\addlegendentry{M}

\addplot[thick] coordinates {(0.10,0) (0.15,1) (0.20,1) (0.30,0)};
\addlegendentry{B}

\addplot[thick] coordinates {(0.25,0) (0.30,1) (0.40,1) (0.50,0)};
\addlegendentry{VB}

\end{axis}
\end{tikzpicture}
\caption{Rozmyta skala prawdopodobieństwa dla zmiennej lingwistycznej \textit{Probability}.}
\label{fig:fuzzy-scale}
\end{figure}

\section{Drzewo niezdatności}

Dla analizowanego systemu skonstruowano drzewo niezdatności ze zdarzeniem szczytowym:

\[
T = \text{brak uruchomienia instalacji gaśniczej podczas pożaru}
\]

Założono, że do zdarzenia szczytowego może doprowadzić wystąpienie któregokolwiek z następujących zdarzeń bazowych:
\begin{itemize}
\item $A$ -- awaria czujnika dymu,
\item $B$ -- awaria centrali sterującej,
\item $C$ -- brak zasilania systemu,
\item $D$ -- awaria instalacji gaśniczej.
\end{itemize}

\section{Oszacowanie prawdopodobieństw zdarzeń bazowych}

Dla poszczególnych zdarzeń bazowych przyjęto następujące oceny lingwistyczne oraz odpowiadające im liczby rozmyte trapezowe:

\begin{center}
\begin{tabular}{cccc}
\toprule
Zdarzenie & Opis & Ocena lingwistyczna & Parametry trapezu \\
\midrule
$A$ & Awaria czujnika dymu & M  & $(0.005,0.01,0.03,0.05)$ \\
$B$ & Awaria centrali sterującej & VM & $(0,0,0.002,0.01)$ \\
$C$ & Brak zasilania systemu & M  & $(0.005,0.01,0.03,0.05)$ \\
$D$ & Awaria instalacji gaśniczej & S  & $(0.03,0.06,0.10,0.15)$ \\
\bottomrule
\end{tabular}
\end{center}

Zatem:
\[
\tilde{A}=(0.005,0.01,0.03,0.05)
\]
\[
\tilde{B}=(0,0,0.002,0.01)
\]
\[
\tilde{C}=(0.005,0.01,0.03,0.05)
\]
\[
\tilde{D}=(0.03,0.06,0.10,0.15)
\]

\section{Wyznaczenie prawdopodobieństwa zdarzenia szczytowego}

Dla bramki OR prawdopodobieństwo zdarzenia szczytowego opisuje wzór:
\[
P(T)=1-(1-\tilde{A})(1-\tilde{B})(1-\tilde{C})(1-\tilde{D})
\]

W obliczeniach zastosowano aproksymację trapezową. Dla liczby rozmytej:
\[
\tilde{X}=(a,b,c,d)
\]
przyjmujemy:
\[
1-\tilde{X}=(1-d,\;1-c,\;1-b,\;1-a)
\]

Natomiast dla dwóch dodatnich liczb trapezowych:
\[
\tilde{X}=(a,b,c,d), \qquad \tilde{Y}=(e,f,g,h)
\]
stosujemy aproksymację:
\[
\tilde{X}\otimes\tilde{Y}\approx(ae,\;bf,\;cg,\;dh)
\]

\subsection{Krok 1 -- wyznaczenie dopełnień}

\[
1-\tilde{A}=(1-0.05,\;1-0.03,\;1-0.01,\;1-0.005)
\]
\[
1-\tilde{A}=(0.95,\;0.97,\;0.99,\;0.995)
\]

\[
1-\tilde{B}=(1-0.01,\;1-0.002,\;1-0,\;1-0)
\]
\[
1-\tilde{B}=(0.99,\;0.998,\;1,\;1)
\]

\[
1-\tilde{C}=(0.95,\;0.97,\;0.99,\;0.995)
\]

\[
1-\tilde{D}=(1-0.15,\;1-0.10,\;1-0.06,\;1-0.03)
\]
\[
1-\tilde{D}=(0.85,\;0.90,\;0.94,\;0.97)
\]

\subsection{Krok 2 -- iloczyn $(1-\tilde{A})(1-\tilde{B})$}

\[
(1-\tilde{A})(1-\tilde{B}) \approx
(0.95\cdot0.99,\;0.97\cdot0.998,\;0.99\cdot1,\;0.995\cdot1)
\]

\[
(1-\tilde{A})(1-\tilde{B}) \approx
(0.9405,\;0.96806,\;0.99,\;0.995)
\]

\subsection{Krok 3 -- iloczyn z $(1-\tilde{C})$}

\[
[(1-\tilde{A})(1-\tilde{B})](1-\tilde{C}) \approx
(0.9405\cdot0.95,\;0.96806\cdot0.97,\;0.99\cdot0.99,\;0.995\cdot0.995)
\]

\[
\approx
(0.893475,\;0.9390182,\;0.9801,\;0.990025)
\]

\subsection{Krok 4 -- iloczyn z $(1-\tilde{D})$}

\[
[(1-\tilde{A})(1-\tilde{B})(1-\tilde{C})](1-\tilde{D}) \approx
\]
\[
(0.893475\cdot0.85,\;0.9390182\cdot0.90,\;0.9801\cdot0.94,\;0.990025\cdot0.97)
\]

\[
\approx
(0.75945375,\;0.84511638,\;0.921294,\;0.96032425)
\]

\subsection{Krok 5 -- wyznaczenie $P(T)$}

Teraz obliczamy dopełnienie otrzymanego wyniku:
\[
P(T)=1-(0.75945375,\;0.84511638,\;0.921294,\;0.96032425)
\]

\[
P(T)\approx
(1-0.96032425,\;1-0.921294,\;1-0.84511638,\;1-0.75945375)
\]

\[
P(T)\approx
(0.03967575,\;0.078706,\;0.15488362,\;0.24054625)
\]

Po zaokrągleniu:
\[
P(T)\approx(0.04,\;0.079,\;0.155,\;0.241)
\]

Dla uproszczenia dalszej interpretacji można przyjąć:
\[
P(T)\approx(0.04,\;0.08,\;0.15,\;0.24)
\]

Oznacza to, że:
\begin{itemize}
\item nośnik liczby rozmytej wynosi w przybliżeniu $[0.04,0.24]$,
\item rdzeń liczby rozmytej wynosi w przybliżeniu $[0.08,0.15]$.
\end{itemize}

\begin{figure}[H]
\centering
\begin{tikzpicture}
\begin{axis}[
    width=12cm,
    height=7cm,
    xmin=0, xmax=0.26,
    ymin=0, ymax=1.1,
    xlabel={$P(T)$},
    ylabel={Stopień przynależności},
    grid=both
]

\addplot[thick] coordinates {(0.0397,0) (0.0787,1) (0.1549,1) (0.2405,0)};

\end{axis}
\end{tikzpicture}
\caption{Aproksymowana trapezowa funkcja przynależności prawdopodobieństwa zdarzenia szczytowego $P(T)$.}
\label{fig:pt-trapezoid}
\end{figure}

\section{Analiza znaczenia zdarzeń bazowych}

Analiza znaczenia polega na sprawdzeniu, jak zmieniają się parametry trapezowej funkcji przynależności $P(T)$ po wyeliminowaniu negatywnego wpływu danego zdarzenia bazowego, czyli po przyjęciu jego prawdopodobieństwa równego:
\[
(0,0,0,0)
\]

\subsection{Eliminacja zdarzenia $A$}

Przyjmujemy:
\[
\tilde{A}=(0,0,0,0)
\]
zatem:
\[
1-\tilde{A}=(1,1,1,1)
\]

Wtedy:
\[
P(T|A=0)=1-(1-\tilde{B})(1-\tilde{C})(1-\tilde{D})
\]

Najpierw:
\[
(1-\tilde{B})(1-\tilde{C})\approx
(0.99\cdot0.95,\;0.998\cdot0.97,\;1\cdot0.99,\;1\cdot0.995)
\]
\[
\approx(0.9405,\;0.96806,\;0.99,\;0.995)
\]

Następnie:
\[
(0.9405,\;0.96806,\;0.99,\;0.995)(0.85,\;0.90,\;0.94,\;0.97)
\]
\[
\approx(0.799425,\;0.871254,\;0.9306,\;0.96515)
\]

Stąd:
\[
P(T|A=0)\approx
(1-0.96515,\;1-0.9306,\;1-0.871254,\;1-0.799425)
\]

\[
P(T|A=0)\approx
(0.03485,\;0.0694,\;0.128746,\;0.200575)
\]

\subsection{Eliminacja zdarzenia $B$}

Przyjmujemy:
\[
\tilde{B}=(0,0,0,0)
\]
zatem:
\[
1-\tilde{B}=(1,1,1,1)
\]

Wtedy:
\[
P(T|B=0)=1-(1-\tilde{A})(1-\tilde{C})(1-\tilde{D})
\]

Najpierw:
\[
(1-\tilde{A})(1-\tilde{C})\approx
(0.95\cdot0.95,\;0.97\cdot0.97,\;0.99\cdot0.99,\;0.995\cdot0.995)
\]
\[
\approx(0.9025,\;0.9409,\;0.9801,\;0.990025)
\]

Następnie:
\[
(0.9025,\;0.9409,\;0.9801,\;0.990025)(0.85,\;0.90,\;0.94,\;0.97)
\]
\[
\approx(0.767125,\;0.84681,\;0.921294,\;0.96032425)
\]

Stąd:
\[
P(T|B=0)\approx
(1-0.96032425,\;1-0.921294,\;1-0.84681,\;1-0.767125)
\]

\[
P(T|B=0)\approx
(0.03967575,\;0.078706,\;0.15319,\;0.232875)
\]

\subsection{Eliminacja zdarzenia $C$}

Ponieważ zdarzenia $A$ i $C$ mają identyczne parametry, otrzymujemy taki sam wynik jak dla eliminacji zdarzenia $A$:

\[
P(T|C=0)\approx
(0.03485,\;0.0694,\;0.128746,\;0.200575)
\]

\subsection{Eliminacja zdarzenia $D$}

Przyjmujemy:
\[
\tilde{D}=(0,0,0,0)
\]
zatem:
\[
1-\tilde{D}=(1,1,1,1)
\]

Wtedy:
\[
P(T|D=0)=1-(1-\tilde{A})(1-\tilde{B})(1-\tilde{C})
\]

Z wcześniejszych obliczeń:
\[
(1-\tilde{A})(1-\tilde{B})(1-\tilde{C})\approx
(0.893475,\;0.9390182,\;0.9801,\;0.990025)
\]

Stąd:
\[
P(T|D=0)\approx
(1-0.990025,\;1-0.9801,\;1-0.9390182,\;1-0.893475)
\]

\[
P(T|D=0)\approx
(0.009975,\;0.0199,\;0.0609818,\;0.106525)
\]

\subsection{Zestawienie wyników}

\begin{center}
\begin{tabular}{cc}
\toprule
Przypadek & Aproksymacja trapezowa $P(T)$ \\
\midrule
Wartość bazowa & $(0.0397,\;0.0787,\;0.1549,\;0.2405)$ \\
bez $A$ & $(0.0349,\;0.0694,\;0.1287,\;0.2006)$ \\
bez $B$ & $(0.0397,\;0.0787,\;0.1532,\;0.2329)$ \\
bez $C$ & $(0.0349,\;0.0694,\;0.1287,\;0.2006)$ \\
bez $D$ & $(0.0100,\;0.0199,\;0.0610,\;0.1065)$ \\
\bottomrule
\end{tabular}
\end{center}

\begin{figure}[H]
\centering
\begin{tikzpicture}
\begin{axis}[
    width=14cm,
    height=8cm,
    xmin=0, xmax=0.26,
    ymin=0, ymax=1.1,
    xlabel={$P(T)$},
    ylabel={Stopień przynależności},
    grid=both,
    legend style={at={(0.5,-0.2)}, anchor=north, legend columns=2}
]

\addplot[thick] coordinates {(0.0397,0) (0.0787,1) (0.1549,1) (0.2405,0)};
\addlegendentry{Wartość bazowa}

\addplot[thick,dashed] coordinates {(0.0349,0) (0.0694,1) (0.1287,1) (0.2006,0)};
\addlegendentry{bez $A$}

\addplot[thick,dotted] coordinates {(0.0397,0) (0.0787,1) (0.1532,1) (0.2329,0)};
\addlegendentry{bez $B$}

\addplot[thick,dashdotted] coordinates {(0.0349,0) (0.0694,1) (0.1287,1) (0.2006,0)};
\addlegendentry{bez $C$}

\addplot[thick] coordinates {(0.0100,0) (0.0199,1) (0.0610,1) (0.1065,0)};
\addlegendentry{bez $D$}

\end{axis}
\end{tikzpicture}
\caption{Porównanie trapezowych funkcji przynależności prawdopodobieństwa zdarzenia szczytowego dla różnych wariantów eliminacji zdarzeń bazowych.}
\label{fig:importance}
\end{figure}

\section{Wnioski}

Na podstawie przeprowadzonej analizy można sformułować następujące wnioski:

\begin{enumerate}
\item Zastosowanie trapezowych liczb rozmytych pozwala uwzględnić niepewność oceny prawdopodobieństw zdarzeń bazowych w bardziej realistyczny sposób niż przy użyciu pojedynczych wartości liczbowych.

\item Prawdopodobieństwo zdarzenia szczytowego zostało wyznaczone w postaci liczby rozmytej trapezowej:
\[
P(T)\approx(0.0397,\;0.0787,\;0.1549,\;0.2405)
\]
co po zaokrągleniu można zapisać jako:
\[
P(T)\approx(0.04,\;0.08,\;0.15,\;0.24)
\]

\item Otrzymany wynik oznacza, że najbardziej wiarygodne wartości prawdopodobieństwa zdarzenia szczytowego znajdują się w przybliżeniu w przedziale od $0.08$ do $0.15$, natomiast pełny zakres możliwych wartości mieści się od około $0.04$ do $0.24$.

\item Analiza znaczenia pokazuje, że największy wpływ na prawdopodobieństwo zdarzenia szczytowego ma zdarzenie $D$, czyli awaria instalacji gaśniczej. Po jego eliminacji wszystkie parametry liczby rozmytej $P(T)$ ulegają największemu zmniejszeniu.

\item Zdarzenia $A$ i $C$ mają identyczny i umiarkowanie duży wpływ na wynik końcowy, co wynika z przyjęcia takich samych parametrów trapezowych dla obu zdarzeń.

\item Najmniejszy wpływ na wynik ma zdarzenie $B$, czyli awaria centrali sterującej. Po jego eliminacji zmiana parametrów funkcji przynależności jest niewielka.

\item W praktyce oznacza to, że działania zwiększające bezpieczeństwo systemu powinny być skierowane przede wszystkim na poprawę niezawodności instalacji gaśniczej, a następnie na zapewnienie poprawnej pracy czujników dymu i zasilania.
\end{enumerate}

\end{document}